% =============================================================
% บทที่ 1 บทนำ — เครื่องมือที่ใช้พัฒนา และขอบเขต/ข้อจำกัดของระบบ
% ข้อมูลตรวจจาก main repo (README.md, CLAUDE.md, package.json) และเว็บจริง dinokhonkaen.app
% =============================================================

\section{เครื่องมือที่ใช้พัฒนา}

เว็บแอปพลิเคชัน ``ไดโน'' (Dino Khon Kaen) เผยแพร่ใช้งานจริงที่ https://dinokhonkaen.app พัฒนาเป็นบริการอิสระ 3 ส่วน ได้แก่ ส่วนติดต่อผู้ใช้งาน บริการข้อมูลหลัก และ AI Service ซึ่งแยกโฟลเดอร์และทำงานแยกจากกัน โดยใช้เครื่องมือดังนี้

\subsection{React และ Vite}

ใช้พัฒนาส่วนติดต่อผู้ใช้งาน (โฟลเดอร์ \texttt{frontend}) ซึ่งเป็นเว็บแอปพลิเคชันชุดเดียวที่ให้บริการทั้งฝั่งผู้ใช้งานและหน้าผู้ดูแลระบบ ใช้ React 18 ร่วมกับ react-router-dom และใช้ Vite เป็นเครื่องมือพัฒนาและสร้างไฟล์สำหรับใช้งานจริง

\subsection{Node.js และ Express}

ใช้พัฒนาบริการข้อมูลหลัก (โฟลเดอร์ \texttt{backend}) ให้บริการ REST API ภายใต้เส้นทาง \texttt{/api} สำหรับจัดการข้อมูลสถานที่ กิจกรรม ฐานความรู้ QR Code ของรางวัล พอยท์ ทริป และระบบสมาชิก พร้อมทั้งเชื่อมต่อฐานข้อมูลผ่านไลบรารี knex

\subsection{Python และ FastAPI}

ใช้พัฒนา AI Service (โฟลเดอร์ \texttt{chatbot-service}) ซึ่งรับผิดชอบแชตบอต ระบบวางแผนทริป การสกัดข้อมูลกิจกรรมจากข้อความโพสต์ และการสร้าง Embedding โดยเรียกโมเดลภาษาขนาดใหญ่ผ่าน KKU AI Gateway

\subsection{PostgreSQL บน Supabase}

ใช้เป็นฐานข้อมูลหลักของระบบ พร้อมส่วนขยาย pgvector สำหรับการค้นคืนข้อมูลเชิงเวกเตอร์ โดยโครงสร้างตารางอยู่ในไฟล์ migration ของโฟลเดอร์ \texttt{backend/supabase/migrations}

\subsection{บริการภายนอก}

ระบบใช้ Cloudinary สำหรับเก็บรูปภาพ Google Places API สำหรับนำเข้าข้อมูลสถานที่ Google Maps JavaScript API สำหรับแสดงแผนที่ และ Gmail SMTP สำหรับส่งอีเมลยืนยันและรีเซ็ตรหัสผ่าน

\subsection{Visual Studio Code}

ใช้เป็นเครื่องมือสำหรับเขียนและแก้ไขซอร์สโค้ดของโปรแกรม

\section{ส่วนประมวลผลปัญญาประดิษฐ์ (AI Service)}

\subsection{ภาพรวมและหน้าที่}

AI Service คือบริการที่พัฒนาด้วย Python และ FastAPI ในโฟลเดอร์ \texttt{chatbot-service} ทำหน้าที่ประมวลผลทุกฟังก์ชันที่ใช้ปัญญาประดิษฐ์ของระบบ โดยเรียกใช้โมเดลภาษาขนาดใหญ่ (Large Language Model: LLM) ผ่าน KKU AI Gateway ซึ่งมีรูปแบบการเรียกใช้เข้ากันได้กับ OpenAI และสามารถเปลี่ยนไปใช้ Vertex AI ของ Google แทนได้ด้วยการตั้งค่า ฟังก์ชันที่ AI Service ให้บริการแสดงดังตารางที่ \ref{tab:ai-functions}

\begin{table}[H]
\caption{ฟังก์ชันของ AI Service}
\label{tab:ai-functions}
\centering
\small
\begin{tabular}{|L{3.0cm}|L{3.3cm}|L{3.0cm}|L{3.1cm}|}
\hline
\textbf{ฟังก์ชัน} & \textbf{เส้นทาง} & \textbf{ผู้เรียกใช้} & \textbf{ใช้ LLM} \\
\hline
แชตบอตน้องไดโน & \texttt{/chat/} & เบราว์เซอร์ของผู้ใช้ & ใช้ (ปรับคำถาม และตอบ) \\
\hline
วางแผนทริป & \texttt{/trip/llm} & เบราว์เซอร์ของผู้ใช้ & ใช้ (จัดแผน และตรวจคุณภาพ) \\
\hline
สกัดข้อมูลกิจกรรมจากโพสต์ & \texttt{/events/extract} & backend (ผู้ดูแลระบบ) & ใช้ \\
\hline
สร้างคำอธิบายสถานที่ & \texttt{/admin/places/describe} & backend (ผู้ดูแลระบบ) & ใช้ \\
\hline
สร้างดัชนีค้นหา (Embedding) & \texttt{/admin/reindex} & backend (ผู้ดูแลระบบ) & ไม่ใช้ \\
\hline
\end{tabular}
\end{table}

เส้นทางสำหรับผู้ดูแลระบบของ AI Service ไม่มีการตรวจสิทธิ์ในตัวเอง จึงถูกกำหนดให้เรียกผ่าน backend ซึ่งเป็นผู้ตรวจสิทธิ์ผู้ดูแลระบบเท่านั้น และไม่เปิดให้เบราว์เซอร์เรียกโดยตรง

\subsection{ข้อมูลและโมเดลที่ใช้}

AI Service ใช้ข้อมูล 3 ชุดในฐานข้อมูล ได้แก่ ข้อมูลสถานที่ ฐานความรู้ของแชตบอต และข้อมูลกิจกรรม โดยแปลงข้อความของแต่ละรายการเป็นเวกเตอร์ (Embedding) ขนาด 384 มิติ ด้วยโมเดล \texttt{paraphrase-multilingual-MiniLM-L12-v2} ซึ่งรองรับภาษาไทยและทำงานบนเครื่องของ AI Service เอง ไม่ต้องเรียกบริการภายนอก จากนั้นเก็บไว้ในคอลัมน์ \texttt{embedding} ของฐานข้อมูล PostgreSQL ที่เปิดใช้ส่วนขยาย pgvector เพื่อใช้ค้นหาตามความหมาย

การค้นหาใช้ 2 วิธีร่วมกัน (hybrid search) คือ ค้นตามความคล้ายของเวกเตอร์ และค้นตามคำสำคัญในชื่อและป้ายกำกับของสถานที่ เพื่อไม่ให้พลาดรายการที่ตอบคำถามได้จริงแต่เวกเตอร์ไม่ใกล้เคียงมากพอ นอกจากนี้ระบบขยายคำถามด้วยคำพ้องความหมายก่อนค้นหา ส่วนกิจกรรมที่ถูกยกเลิกหรือสิ้นสุดแล้วจะไม่ถูกนำมาตอบ

เมื่อผู้ดูแลระบบเพิ่มหรือแก้ไขข้อมูล รายการนั้นจะถูกทำเครื่องหมายว่ารอสร้างเวกเตอร์ใหม่ และต้องกดปุ่ม ``อัปเดตดัชนีค้นหา'' ในหน้าแดชบอร์ดผู้ดูแลระบบ เพื่อให้ระบบสร้างเวกเตอร์เฉพาะรายการที่รออยู่ การสร้างเวกเตอร์ใหม่ทั้งหมดใช้สคริปต์ \texttt{scripts/embed\_content.py} และต้องทำเมื่อเปลี่ยนโมเดล Embedding

\subsection{การทำงานของแชตบอตน้องไดโน}

แชตบอตใช้เทคนิค Retrieval-Augmented Generation (RAG) คือ ค้นข้อมูลที่เกี่ยวข้องจากฐานข้อมูลก่อน แล้วให้ LLM เรียบเรียงคำตอบโดยอ้างอิงเฉพาะข้อมูลที่ค้นได้ ขั้นตอนการทำงานในแต่ละคำถามมีดังนี้

\begin{numlist}
\item รับข้อความของผู้ใช้ (ไม่เกิน 500 ตัวอักษร) พร้อมประวัติสนทนาล่าสุดไม่เกิน 4 ข้อความ ซึ่งฝั่งผู้ใช้เป็นผู้เก็บประวัติ ส่วน AI Service ไม่เก็บสถานะการสนทนา
\item หากมีประวัติสนทนา ระบบให้ LLM เรียบเรียงคำถามใหม่ให้สมบูรณ์ในตัวเอง เช่น คำถามต่อเนื่องว่า ``แล้วร้านที่สองล่ะ'' จะถูกขยายเป็นคำถามที่ระบุชื่อสถานที่ชัดเจน เพื่อให้ค้นหาได้ถูกต้อง
\item ค้นข้อมูลที่เกี่ยวข้องอย่างละไม่เกิน 3 รายการ จากสถานที่ ฐานความรู้ และกิจกรรม
\item ส่งข้อมูลที่ค้นได้และประวัติสนทนาให้ LLM ตอบ โดยกำหนดให้ตอบจากข้อมูลที่ได้รับเท่านั้น และให้ระบุแหล่งที่มาของคำตอบกำกับไว้ในรูปแบบแท็ก
\item ระบบอ่านแท็กดังกล่าวเพื่อตัดสินว่าคำตอบอ้างอิงข้อมูลชุดใด แล้วแนบการ์ดสถานที่หรือกิจกรรมที่ใช้อ้างอิงไปกับคำตอบ หากไม่พบข้อมูลที่เกี่ยวข้อง แชตบอตจะตอบข้อความแจ้งว่าไม่มีข้อมูลในส่วนนี้ แทนการคาดเดา
\item ส่งคำตอบกลับแบบสตรีม (Server-Sent Events) เป็นช่วง ๆ เพื่อให้ผู้ใช้เห็นคำตอบทยอยแสดงโดยไม่ต้องรอจนเสร็จ
\end{numlist}

\subsection{การทำงานของระบบวางแผนทริป}

ระบบวางแผนทริปเป็นแบบผสม (hybrid) คือ ให้ LLM ทำหน้าที่เลือกและจัดลำดับสถานที่ ส่วนการคำนวณระยะทาง เวลาเดินทาง เวลาเปิด-ปิด และค่าใช้จ่ายใช้โปรแกรมที่คำนวณได้แน่นอน (deterministic) เพื่อไม่ให้ผลลัพธ์ที่เป็นตัวเลขผิดพลาดจากการสร้างข้อความของโมเดล ขั้นตอนการทำงานมีดังนี้

\begin{numlist}
\item รับข้อมูลจากแบบฟอร์ม 4 ขั้นตอน ได้แก่ จำนวนวันและวันเริ่มต้น ที่พัก สถานที่ที่ต้องไป ความสนใจ งบประมาณ จังหวะการเที่ยว ขอบเขตพื้นที่ และช่วงเวลาของแต่ละวัน
\item สร้างรายการสถานที่ตัวเลือก โดยค้นตามความสนใจแต่ละข้อ กรองตามระดับงบประมาณและขอบเขตพื้นที่ และเผื่อจำนวนตัวเลือกให้มากกว่าที่ต้องใช้ เพื่อให้ระบบมีสถานที่สำรองเติมช่วงเวลาที่ว่าง
\item ให้ LLM เลือกและจัดสถานที่เป็นรายวัน โดยกำหนดให้ตอบเป็น JSON และมีข้อจำกัด เช่น ร้านอาหารไม่เกิน 2 แห่งต่อวัน และไม่นำคาเฟ่มาเรียงต่อกันหลายแห่ง
\item โปรแกรมคำนวณตารางเวลาจริง ได้แก่ ลำดับการเดินทาง เวลาถึงและเวลาออก ระยะทางรวม และค่าใช้จ่ายโดยประมาณ พร้อมตรวจเวลาเปิด-ปิดของสถานที่ และเติมสถานที่สำรองหากวันใดยังว่างอยู่
\item ให้ LLM อีกชุดทำหน้าที่ผู้ตรวจ (judge) ประเมินว่าแผนเหมาะกับจังหวะการเที่ยวและตรงกับความสนใจของผู้ใช้หรือไม่ หากไม่ผ่าน ระบบจะสร้างแผนใหม่พร้อมข้อเสนอแนะ โดยทำซ้ำได้ไม่เกิน 2 รอบ
\item ส่งผลลัพธ์กลับพร้อมสรุปเหตุผลของการจัดแผน และแจ้งข้อควรระวัง เช่น สถานที่ที่ระบุแต่จัดเข้าแผนไม่ได้ หรือช่วงเดินทางที่ตรงกับวันหยุดนักขัตฤกษ์
\end{numlist}

หากโมเดลตอบรูปแบบไม่ถูกต้องหรือการวางแผนล้มเหลว ระบบจะส่งข้อความ ``วางแผนทริปไม่สำเร็จ กรุณาลองใหม่อีกครั้ง'' กลับไป โดยไม่สร้างแผนสำรองด้วยการคาดเดา

\subsection{การสกัดข้อมูลกิจกรรมและการสร้างคำอธิบายสถานที่}

ฟังก์ชันสำหรับผู้ดูแลระบบ 2 ฟังก์ชันใช้ LLM ช่วยลดงานกรอกข้อมูล โดยผลลัพธ์ทุกครั้งต้องให้ผู้ดูแลระบบตรวจสอบก่อนบันทึก

\begin{parenlist}
\item การสกัดข้อมูลกิจกรรม รับข้อความโพสต์ภาษาไทยแล้วให้ LLM แยกเป็นช่องข้อมูลของแบบฟอร์มกิจกรรม เช่น ชื่องาน ช่วงวันที่ สถานที่ ค่าเข้าชม และผู้จัดงาน โดยกำหนดให้ใช้เฉพาะข้อมูลที่ปรากฏในโพสต์ หากไม่ระบุให้เว้นว่าง และระบบตรวจรูปแบบวันที่ซ้ำ ตัดช่วงวันที่ที่ผิดพลาดออก
\item การสร้างคำอธิบายสถานที่ นำชื่อ หมวดหมู่ ป้ายกำกับ และรีวิวจาก Google มาสรุปเป็นคำอธิบาย และหากกำหนดคีย์ Serper ไว้ จะใช้ผลค้นหาจากเว็บประกอบด้วย หากข้อมูลไม่พอระบบใช้ข้อความแม่แบบแทน และสามารถสลับไปใช้โมเดลสำรองเมื่อโควตารายวันของโมเดลหลักหมด
\end{parenlist}

\subsection{การจำกัดการใช้งานและการติดตามค่าใช้จ่าย}

เนื่องจากการเรียกใช้ LLM มีค่าใช้จ่ายตามปริมาณโทเคน AI Service จึงมีมาตรการควบคุม ดังนี้

\begin{parenlist}
\item จำกัดจำนวนคำขอต่อผู้ใช้ (ตาม IP) โดยค่าเริ่มต้นคือ แชตบอต 20 ครั้งต่อนาทีและ 200 ครั้งต่อชั่วโมง และการวางแผนทริป 5 ครั้งต่อ 10 นาทีและ 20 ครั้งต่อชั่วโมง
\item จำกัดจำนวนคำขอรวมทุกผู้ใช้ ค่าเริ่มต้นคือ แชตบอต 120 ครั้งต่อนาทีและ 3,000 ครั้งต่อชั่วโมง และการวางแผนทริป 30 ครั้งต่อ 10 นาทีและ 200 ครั้งต่อชั่วโมง เพื่อกำหนดเพดานค่าใช้จ่ายสูงสุด เมื่อเกินจะตอบกลับสถานะ 429
\item จำกัดความยาวข้อความที่รับเข้า เช่น ข้อความแชตไม่เกิน 500 ตัวอักษร และจำนวนวันของทริปไม่เกิน 14 วัน
\item บันทึกจำนวนโทเคนและค่าใช้จ่ายโดยประมาณของทุกคำขอลงไฟล์ \texttt{logs/usage.jsonl} และสรุปผลได้ด้วยสคริปต์ \texttt{scripts/usage\_report.py}
\end{parenlist}

ค่าจำกัดข้างต้นปรับได้ด้วยตัวแปรสภาพแวดล้อมชุด \texttt{RATE\_LIMIT\_*} ในไฟล์ \texttt{chatbot.env} ตามที่ระบุไว้ในไฟล์ \texttt{.env.example}

\section{ขอบเขตและข้อจำกัด}

\subsection{ขอบเขตของผู้ใช้งานทั่วไป (Guest)}

\begin{parenlist}
\item เข้าชมหน้าเว็บสาธารณะ ค้นหาและดูรายละเอียดสถานที่ท่องเที่ยว ร้านอาหาร คาเฟ่ และกิจกรรมได้
\item สอบถามข้อมูลการท่องเที่ยวกับน้องไดโน (แชตบอต) ได้
\item สร้างแผนการเดินทางด้วย AI ได้ แต่แผนจะสูญหายเมื่อปิดหน้าเว็บ เนื่องจากไม่ได้บันทึกลงบัญชี
\item ดูรายการสถานที่ร่วมรายการและของรางวัลในหน้าพอยท์สะสมได้ แต่สะสมพอยท์ไม่ได้
\end{parenlist}

\subsection{ขอบเขตของสมาชิก (Tourist)}

\begin{parenlist}
\item สมัครสมาชิกด้วยอีเมล ยืนยันอีเมล และเข้าสู่ระบบ
\item ใช้งานได้ทุกฟังก์ชันของผู้ใช้งานทั่วไป และบันทึกแผนทริปไว้ในเมนู ``ทริปของฉัน''
\item สแกน QR Code ณ สถานที่ร่วมรายการเพื่อสะสมพอยท์ และดูประวัติพอยท์
\item แลกของรางวัลที่เคาน์เตอร์โดยแจ้งเจ้าหน้าที่
\item แก้ไขข้อมูลโปรไฟล์ เปลี่ยนรหัสผ่านหรืออีเมล ส่งกิจกรรมใหม่เข้าระบบ และแจ้งข้อมูลสถานที่หรือกิจกรรมที่ไม่ถูกต้อง
\end{parenlist}

\subsection{ขอบเขตของผู้ดูแลระบบ (Admin)}

\begin{parenlist}
\item เข้าสู่ระบบผ่านหน้าเข้าสู่ระบบสำหรับผู้ดูแลระบบ
\item ดูภาพรวมของระบบผ่านแดชบอร์ด
\item เพิ่ม แก้ไข และลบข้อมูลสถานที่ กิจกรรม และฐานความรู้ของแชตบอต รวมถึงปรับปรุงข้อมูลสถานที่ให้ตรงกับ Google Places
\item ตรวจสอบรายงานข้อมูลที่ไม่ถูกต้องและคำขอกิจกรรมจากผู้ใช้
\item จัดการ QR Code ของรางวัล และดำเนินการแลกของรางวัลที่เคาน์เตอร์
\item จัดการผู้ใช้ ดูทริปที่ผู้ใช้สร้าง และดูผลความพึงพอใจ
\end{parenlist}

\subsection{ข้อจำกัดของระบบ}

\begin{parenlist}
\item ข้อมูลสถานที่ครอบคลุมจังหวัดขอนแก่น โดยข้อมูลหลักมาจาก Google Places API และข้อมูลที่ผู้ดูแลระบบเพิ่มเอง จึงอาจไม่เป็นปัจจุบันหากยังไม่ได้ปรับปรุงข้อมูล
\item แชตบอตและระบบวางแผนทริปใช้โมเดลภาษาขนาดใหญ่ แม้จำกัดการตอบให้อ้างอิงฐานความรู้ แต่ยังมีโอกาสตอบคลาดเคลื่อน และอาจใช้งานไม่ได้ชั่วคราวเมื่อโควตาของ KKU AI Gateway หมด
\item QR Code แต่ละจุดให้พอยท์ได้ครั้งเดียวต่อบัญชี และการแลกของรางวัลต้องทำที่เคาน์เตอร์โดยเจ้าหน้าที่ ระบบยังไม่รองรับการแลกของรางวัลด้วยตนเอง
\item ผู้ดูแลระบบเป็นบัญชีที่สร้างโดยนักพัฒนาผ่านสคริปต์ ไม่สามารถสมัครผ่านหน้าเว็บไซต์ได้
\item ระบบเป็นเว็บแอปพลิเคชันที่ใช้ผ่านเบราว์เซอร์ ไม่มีแอปพลิเคชันบนมือถือแยกต่างหาก และการสแกน QR Code ต้องอนุญาตให้เว็บเข้าถึงกล้อง
\end{parenlist}
